\documentclass{article}

\usepackage[utf8]{inputenc}
\usepackage[T1]{fontenc}
\usepackage{amsmath,amssymb,amsthm}
\usepackage{mathtools}
\usepackage{hyperref}
\usepackage{booktabs}

\makeatletter
\newcommand{\citep}[2][]{%
  \begingroup
  \def\@tmpkey{#2}%
  \ifx\@tmpkey\@empty\else
    (\ifx\@empty#1\@empty\else #1, \fi
     \@nameuse{bibkey@\@tmpkey})%
  \fi
  \endgroup
}
\expandafter\def\csname bibkey@lasser2026phase\endcsname{Phase~Redundancy}
\expandafter\def\csname bibkey@lasser2026micro\endcsname{Microfoundation}
\makeatother

\newtheorem{definition}{Definition}
\newtheorem{lemma}{Lemma}
\newtheorem{theorem}{Theorem}
\newtheorem{remark}{Remark}
\newtheorem{assumption}{Assumption}

\DeclareMathOperator{\vol}{vol}
\newcommand{\volL}{{\vol_{\mathrm{P}}}}
\newcommand{\volR}{{\vol_{\mathrm{R}}}}
\newcommand{\volRwin}{{\vol_{\mathrm{R}}^{[W]}}}
\newcommand{\Pproxy}{P}
\newcommand{\Ttruth}{T}
\newcommand{\Lyap}{L}
\newcommand{\Wm}{\hat{W}}
\newcommand{\epsgap}{\varepsilon_{\mathrm{gap}}}
\newcommand{\epsnonres}{\varepsilon_{\mathrm{gap}}^{\mathrm{nonres}}}
\newcommand{\epsfloor}{\varepsilon_{\mathrm{floor}}^{\mathrm{res}}}

\title{Lemma 2 (Lyapunov-Goodhart bridge): Full Proof Draft}
\author{Paper 10 working draft for codex review}
\date{}

\begin{document}

\maketitle

\section{Goal of this draft}

Promote Lemma~2 of Paper~10 \S4.3 from inline proof outline to a
self-contained formal proof, suitable for placement in the
appendix.

The current outline asserts a Cauchy-Schwarz argument relating the
Lyapunov function $\Lyap = \sum_k w_k \epsilon_k^2$ from
\citep[Phase~Redundancy]{lasser2026phase} to the proxy-truth sup-norm
$\epsnonres = \sup_{c} |T(c) - P(c)|$ from
\citep[Microfoundation]{lasser2026micro}, but does not fully
discharge the per-capability dimension structure or address the
coordinate-Lipschitz parameterization assumption required for the
per-capability gap bound.

This draft formalizes the proof, names two deployment-class
conditions explicitly (bounded per-capability dimension count;
coordinate-Lipschitz parameterization), and walks through the
Cauchy-Schwarz step-by-step.

\section{Deployment-class conditions}

The proof uses three deployment-class conditions, each
characterizing the world-model parameterization choice. Each is a
property of the deployment class $\mathcal{D}$, fixed before $|P|$
is varied (per Definition~1 of Lemma~1).

\begin{assumption}[BD: Bounded per-capability dimension count]
\label{ass:BD}
There exists a constant $N_{\max}$ depending on $\mathcal{D}$ but
independent of $|P|$ such that, for every capability $c \in
P^{\mathrm{act}}$ in the operationally active subspace, the world-
model representation of $c$ involves at most $N_{\max}$ world-model
dimensions:
\[
  |\mathrm{dim}(c)| \;\leq\; N_{\max}, \qquad
  \mathrm{dim}(c) := \{k \in \{1, \ldots, K\} :
  \text{dim. } k \text{ enters } P(c) \text{ or } T(c)\}.
\]
\end{assumption}

\textbf{Operational interpretation.}  Standard world-model
parameterizations describe each capability with a finite, bounded
set of properties (existence indicator, attestation status,
exercise frequency, leverage estimate, etc.). The number of
relevant properties per capability does not scale with the total
number of capabilities. Pathological parameterizations where each
capability adds $|P|$-many dimensions (e.g., pairwise-interaction
parameterizations, where $\mathrm{dim}(c)$ scales with $|P|$ via
all-pairs cross-effects) violate BD and are excluded from
$\mathcal{D}$.

\begin{assumption}[CL: Coordinate-Lipschitz parameterization]
\label{ass:CL}
For every capability $c \in P^{\mathrm{act}}$, the per-capability
proxy-truth gap is bounded by the sum of per-dimension errors over
$\mathrm{dim}(c)$:
\[
  |T(c) - P(c)| \;\leq\; \sum_{k \in \mathrm{dim}(c)} |\epsilon_k|,
  \qquad \epsilon_k := \Wm_k - \Wm_k^*.
\]
where $\epsilon_k$ is the per-dimension world-model error from
\citep[Phase~Redundancy]{lasser2026phase}.

Equivalently: the parameterization is coordinate-Lipschitz with
constant $\leq 1$ in each dimension. (Coordinate-Lipschitz with
constant $L_k$ in dimension $k$ would replace $|\epsilon_k|$ by
$L_k \cdot |\epsilon_k|$ in the sum; we set $L_k = 1$ by absorbing
the $L_k$ factor into the world-model weighting $w_k$ without loss
of generality.)
\end{assumption}

\textbf{Operational interpretation.}  CL is a regularity property
of the world-model parameterization. If the world-model represents
the proxy and truth values for each capability as a smooth function
of the parameter vector $\Wm$, with bounded coordinate-wise
sensitivity, then CL holds. Pathological parameterizations where
small parameter perturbations produce arbitrary capability-level
gaps (e.g., bilinear terms, threshold non-linearities) violate
CL and are excluded from $\mathcal{D}$.

\begin{assumption}[WB: Bounded weight floor]
\label{ass:WB}
The per-dimension weights $w_k$ in
\citep[Phase~Redundancy]{lasser2026phase}'s Lyapunov function
$\Lyap(\Wm) = \sum_k w_k \epsilon_k(\Wm)^2$ are uniformly bounded
below by some $w_{\min} > 0$ depending on $\mathcal{D}$ but
independent of $|P|$:
\[
  \forall k \in \{1, \ldots, K\}, \quad w_k \geq w_{\min}.
\]
\end{assumption}

\textbf{Operational interpretation.}  The Lyapunov weighting
encodes safety-relevance per dimension; WB requires every dimension
to have at least minimal safety-relevance. Dimensions with
arbitrarily small weight are absorbed into noise; if the deployment
parameterization includes truly negligible dimensions, those should
be removed before applying the Lyapunov machinery.

\section{Statement of Lemma 2}

\begin{lemma}[Lyapunov-Goodhart quantitative bridge]
\label{lem:2}
Under conditions BD (bounded per-capability dimension count), CL
(coordinate-Lipschitz parameterization), and WB (bounded weight
floor) of the deployment class $\mathcal{D}$, and under invariants
$I_1$ and $I_4$ (which together imply $I_2$:
$\Lyap < \epsilon_{\mathrm{safe}}$ via
\citep[Phase~Redundancy]{lasser2026phase}'s contraction analysis),
the proxy-truth sup-norm gap satisfies:
\[
  \Lyap(\Wm_t) \;<\; \epsilon_{\mathrm{safe}}
  \quad \implies \quad
  \epsnonres \;<\; f(\epsilon_{\mathrm{safe}})
  \;:=\; \sqrt{N_{\max} \cdot \epsilon_{\mathrm{safe}} / w_{\min}}.
\]
The function $f$ is intensive in $|P|$ over $\mathcal{D}$: $f$
depends on $\mathcal{D}$'s operational parameters
($N_{\max}, w_{\min}$) but not on $|P|$.
\end{lemma}

\section{Proof}

\begin{proof}
We work over a deployment class $\mathcal{D}$ satisfying BD, CL,
and WB throughout.

\emph{Step 1 (per-capability gap bound).}  By Assumption CL, for
every capability $c \in P^{\mathrm{act}}$:
\[
  |T(c) - P(c)| \;\leq\; \sum_{k \in \mathrm{dim}(c)} |\epsilon_k|.
\]
This bounds the per-capability proxy-truth gap by the sum of
per-dimension errors over the dimensions that represent $c$.

\emph{Step 2 (Cauchy-Schwarz on per-capability sum).}  Apply
Cauchy-Schwarz to the sum $\sum_{k \in \mathrm{dim}(c)} |\epsilon_k|$
with weights $a_k = \mathbb{1}[k \in \mathrm{dim}(c)]$ and
$b_k = |\epsilon_k|$:
\[
  \sum_{k=1}^{K} a_k b_k \;\leq\; \sqrt{\sum_{k=1}^{K} a_k^2}
  \cdot \sqrt{\sum_{k=1}^{K} b_k^2}.
\]
The first factor:
\[
  \sqrt{\sum_{k=1}^{K} a_k^2} \;=\; \sqrt{|\mathrm{dim}(c)|}
  \;\leq\; \sqrt{N_{\max}}
\]
by Assumption BD. The second factor:
\[
  \sqrt{\sum_{k=1}^{K} b_k^2} \;=\; \sqrt{\sum_{k=1}^{K} \epsilon_k^2}.
\]
Combining:
\[
  \sum_{k \in \mathrm{dim}(c)} |\epsilon_k|
  \;\leq\; \sqrt{N_{\max}} \cdot \sqrt{\sum_{k=1}^{K} \epsilon_k^2}.
\]

\emph{Step 3 (relate sum-of-squares to Lyapunov).}  By Assumption
WB, $w_k \geq w_{\min} > 0$ for all $k$. Therefore:
\[
  \sum_{k=1}^{K} \epsilon_k^2 \;=\; \sum_{k=1}^{K} \frac{w_k}{w_k} \epsilon_k^2
  \;\leq\; \sum_{k=1}^{K} \frac{w_k}{w_{\min}} \epsilon_k^2
  \;=\; \frac{1}{w_{\min}} \sum_{k=1}^{K} w_k \epsilon_k^2
  \;=\; \frac{\Lyap(\Wm)}{w_{\min}}.
\]

\emph{Step 4 (combine into per-capability bound).}  Substituting
Step~3 into Step~2:
\[
  \sum_{k \in \mathrm{dim}(c)} |\epsilon_k|
  \;\leq\; \sqrt{N_{\max}} \cdot \sqrt{\Lyap(\Wm) / w_{\min}}
  \;=\; \sqrt{N_{\max} \cdot \Lyap(\Wm) / w_{\min}}.
\]
Combining with Step~1:
\[
  |T(c) - P(c)| \;\leq\; \sqrt{N_{\max} \cdot \Lyap(\Wm) / w_{\min}}.
\]
The right-hand side is independent of $c$.

\emph{Step 5 (sup over $P^{\mathrm{act}}$).}  Taking supremum over
$c \in P^{\mathrm{act}}$:
\[
  \epsnonres \;=\; \sup_{c \in P^{\mathrm{act}}} |T(c) - P(c)|
  \;\leq\; \sqrt{N_{\max} \cdot \Lyap(\Wm) / w_{\min}}.
\]

\emph{Step 6 (apply Lyapunov bound).}  Under invariants $I_1$ and
$I_4$, \citep[Phase~Redundancy]{lasser2026phase}'s contraction
analysis (Theorem~1a of Phase Redundancy) gives
$\Lyap(\Wm_t) < \epsilon_{\mathrm{safe}}$ in the self-correcting
basin. Substituting:
\[
  \epsnonres \;<\; \sqrt{N_{\max} \cdot \epsilon_{\mathrm{safe}} / w_{\min}}
  \;=\; f(\epsilon_{\mathrm{safe}}).
\]

\emph{Step 7 (intensivity over $\mathcal{D}$).}  The bound
$f(\epsilon_{\mathrm{safe}}) = \sqrt{N_{\max} \cdot \epsilon_{\mathrm{safe}} /
w_{\min}}$ depends on:
\begin{itemize}
\item $N_{\max}$: deployment-class constant by Assumption BD,
  independent of $|P|$.
\item $\epsilon_{\mathrm{safe}}$: deployment-class threshold from
  \citep[Phase~Redundancy]{lasser2026phase}'s self-correcting
  basin, depending on the dynamics' constants
  (correction rate, error-proportional rate, drift rate, subsumption frequency) which are operational rate
  parameters, not $|P|$-dependent quantities.
\item $w_{\min}$: deployment-class constant by Assumption WB,
  independent of $|P|$.
\end{itemize}
Therefore $f(\epsilon_{\mathrm{safe}})$ is intensive in $|P|$ over
$\mathcal{D}$, with bounding constant determined entirely by
deployment-class parameters. $\Box$
\end{proof}

\section{Discussion of constants}

The proof produces $f(\epsilon_{\mathrm{safe}}) = \sqrt{N_{\max}
\cdot \epsilon_{\mathrm{safe}} / w_{\min}}$. The constants:
\begin{itemize}
\item $N_{\max}$: bounded per-capability dimension count
  (\emph{deployment-policy-derived}, BD).  Operationally: a
  property of the world-model parameterization choice. Calibrated
  by enumerating the dimensions used per-capability in the
  deployment's parameter vector.
\item $\epsilon_{\mathrm{safe}}$: Phase Redundancy's safety
  threshold (\emph{source-derived}, from
  \citep[Phase~Redundancy]{lasser2026phase}'s policy-correctness
  analysis).
\item $w_{\min}$: minimum per-dimension weight
  (\emph{deployment-policy-derived}, WB).  Set by the deployment's
  safety-relevance weighting; calibrated by inspecting the
  Lyapunov-weight assignment.
\end{itemize}

The bridge is square-root in $\epsilon_{\mathrm{safe}}$, which
reflects the underlying $L^2 \to L^\infty$ Cauchy-Schwarz
relationship between the weighted-Euclidean Lyapunov norm and the
sup-norm proxy-truth gap. Operators can tighten the bound by
either reducing $N_{\max}$ (simpler parameterization) or increasing
$w_{\min}$ (uniform safety-relevance weighting), at deployment-
engineering cost.

\section{Connection to Theorem 1 Layer 1}

Lemma~2's role in
Theorem~\ref{thm:deployment_safety}'s Layer~1 proof:
\begin{itemize}
\item Steps 2--3 of the Layer~1 proof use Phase Redundancy's
  contraction to establish $\Lyap < \epsilon_{\mathrm{safe}}$
  under invariants $I_1, I_4$.
\item Step~3 invokes Lemma~\ref{lem:2} to convert this Lyapunov
  bound into a sup-norm gap bound:
  $\epsnonres < f(\epsilon_{\mathrm{safe}})$.
\item Step~5 then composes via Microfoundation's Goodhart bound:
  $|g(T) - g(P)| \leq \mathrm{Lip}(g) \cdot \epsnonres <
  \mathrm{Lip}(g) \cdot f(\epsilon_{\mathrm{safe}})$.
\item Step~6 (now via Lemma~\ref{lem:1}'s generic-X form) packages
  this with the residual floor and Lipschitz multiplication, giving
  the intensive Layer~1 bound.
\end{itemize}

Lemma~2 is therefore the bridge that converts world-model accuracy
(Phase Redundancy's deliverable) into capability-level gap accuracy
(Microfoundation's input). Without Lemma~2, the Phase Redundancy
machinery would not connect to the Goodhart bound; without Phase
Redundancy, the Lyapunov bound on $\epsnonres$ would not have an
operational invariant $I_2$ supporting it.

\section{What this proof does and does not establish}

\textbf{Establishes:}
\begin{itemize}
\item A formal Cauchy-Schwarz derivation bridging
  \citep[Phase~Redundancy]{lasser2026phase}'s Lyapunov function
  to \citep[Microfoundation]{lasser2026micro}'s sup-norm gap.
\item Three deployment-class conditions (BD, CL, WB) sufficient
  for the bridge to hold, each with operational interpretation.
\item An explicit intensive-in-$|P|$ bound $f(\epsilon_{\mathrm{safe}})
  = \sqrt{N_{\max} \cdot \epsilon_{\mathrm{safe}} / w_{\min}}$ in
  terms of deployment-class constants.
\item The connection to Theorem 1 Layer 1's intensivity argument,
  via Lemma 1's generic-X composition.
\end{itemize}

\textbf{Does not establish (deferred):}
\begin{itemize}
\item Tightness of $f$. The square-root form is the
  Cauchy-Schwarz upper bound; deployments may achieve sharper
  bounds with additional structure (e.g., per-capability errors
  that are not adversarially aligned).
\item Existence of deployment configurations satisfying BD, CL, WB.
  These are conditions on the parameterization choice; verification
  is by inspection of the deployment's world-model design.
\item Numerical values of $N_{\max}, w_{\min},
  \epsilon_{\mathrm{safe}}$. These are deployment-policy-derived
  constants whose calibration is part of the deployment-tooling
  specification.
\end{itemize}

\section{Specific verification questions for codex review}

\begin{enumerate}
\item Is Assumption CL (coordinate-Lipschitz parameterization)
  correctly stated as the precondition for Step~1's per-capability
  gap bound? In particular: does $|T(c) - P(c)| \leq \sum_{k \in
  \mathrm{dim}(c)} |\epsilon_k|$ require the coordinate-Lipschitz
  constant to equal 1 (as in the absorption argument), or should
  the lemma admit constants $L_k$ explicitly?

\item Is Assumption WB (bounded weight floor) defensible? In
  particular: does
  \citep[Phase~Redundancy]{lasser2026phase}'s Lyapunov machinery
  require $w_k \geq w_{\min} > 0$, or can it accommodate $w_k = 0$
  for some dimensions (with corresponding restrictions)?

\item Is the Cauchy-Schwarz step (Step~2) rigorous? In particular,
  is the indicator-function $a_k = \mathbb{1}[k \in
  \mathrm{dim}(c)]$ choice the right way to bound the per-capability
  sum, or is there a tighter bound using the actual coordinate
  sensitivities?

\item Is Step~3 (sum-of-squares to Lyapunov) rigorous?
  Specifically, does $\sum \epsilon_k^2 \leq \Lyap / w_{\min}$
  preserve all the structure we need, or does it lose information
  that a tighter bound (e.g., using $w_k$ directly) would retain?

\item Is the Layer 1 connection (\S6 of this draft) correct?
  Specifically, does the sequence
  $\Lyap < \epsilon_{\mathrm{safe}} \to \epsnonres < f \to
  |g(T) - g(P)| < \mathrm{Lip}(g) \cdot f$ correctly trace the
  Step 3, 5, 6 of Theorem 1's Layer 1 proof?
\end{enumerate}

\end{document}
